\subsection{Fundamentação Biológica}
\label{subsec:fundamentacao-biologica}

A evolução biológica ocorre em populações de indivíduos que apresentam diferenças hereditárias entre si. De acordo com \citeonline{futuyma2023}, essas diferenças constituem a base da variação evolutiva, pois podem influenciar a sobrevivência e a reprodução dos organismos em determinado ambiente. Ao longo das gerações, características associadas a maior sucesso reprodutivo tendem a se tornar mais frequentes na população.

Segundo o renomado autor \citeonline{darwin1859}, a teoria da evolução por seleção natural constitui a base para a compreensão desse processo, ao demonstrar que indivíduos com características mais favoráveis às condições ambientais tendem a deixar mais descendentes. Posteriormente, os avanços da genética permitiram explicar como características são transmitidas, recombinadas e modificadas entre gerações. Sob o olhar de \citeonline{pierce2024}, conceitos como herança genética, genótipo, fenótipo, recombinação e mutação são fundamentais para compreender a transmissão e a variação das características biológicas. Assim, população, indivíduo, variabilidade genética, herança, genótipo, fenótipo, recombinação, mutação e seleção natural formam a base conceitual necessária para compreender os princípios biológicos que inspiram os Algoritmos Genéticos.

\subsubsection{População, Indivíduo e Ambiente}
\label{subsubsec:populacao-individuo-ambiente}

Na biologia, uma população é formada por indivíduos de uma mesma espécie que compartilham determinado ambiente e participam de processos reprodutivos. Como afrima \citeonline{futuyma2023}, a evolução não ocorre em indivíduos isolados, mas em populações ao longo de gerações, por meio da alteração na frequência de características hereditárias.

Cada indivíduo possui características próprias, algumas herdadas geneticamente, que podem influenciar sua adaptação ao ambiente. Com isso \citeonline{urry2021} diz que, o ambiente exerce influência direta sobre a sobrevivência e a reprodução dos organismos, atuando como fonte de pressão seletiva. Essa adaptação não deve ser entendida como superioridade absoluta, mas como adequação relativa a condições específicas.

\subsubsection{Variabilidade Genética}
\label{subsubsec:variabilidade-genetica}

A variabilidade genética corresponde às diferenças hereditárias existentes entre os indivíduos de uma população. Segundo \citeonline{futuyma2023}, essa variabilidade é essencial para a evolução, pois fornece o conjunto de diferenças sobre o qual a seleção natural pode atuar.

Sem variabilidade, uma população teria menor capacidade de responder a mudanças ambientais. De acordo com\citeonline{urry2021}, a diversidade genética permite que diferentes combinações de características existam em uma mesma população, aumentando as possibilidades de adaptação diante de pressões seletivas. Dessa forma, quanto maior a variabilidade disponível, maior tende a ser o potencial adaptativo da população.

\subsubsection{Herança Genética, Genótipo e Fenótipo}
\label{subsubsec:heranca-genetica-genotipo-fenotipo}

A herança genética é o processo pelo qual características são transmitidas de uma geração para outra. Na visão de \citeonline{pierce2024}, os descendentes recebem material genético de seus progenitores, mas essa transmissão não produz cópias perfeitamente idênticas, pois processos como recombinação e mutação podem alterar as combinações herdadas.

O genótipo corresponde ao conjunto de informações genéticas de um indivíduo. O fenótipo, por sua vez, corresponde às características observáveis ou expressas, resultantes da interação entre genótipo e ambiente. Conforme \citeonline{pierce2024}, essa distinção é importante porque características semelhantes podem surgir de genótipos diferentes, assim como um mesmo genótipo pode produzir fenótipos distintos em ambientes diferentes.

\subsubsection{Seleção Natural}
\label{subsubsec:selecao-natural}

A seleção natural é o processo pelo qual indivíduos com características mais adequadas ao ambiente tendem a sobreviver e se reproduzir com maior frequência. Segundo \citeonline{darwin1859}, esse mecanismo explica como determinadas características podem se tornar mais comuns nas gerações seguintes, quando favorecem a sobrevivência ou a reprodução dos indivíduos que as possuem.

De modo semelhante, \citeonline{futuyma2023} afirma que a seleção natural depende da existência de variação entre indivíduos, da transmissão hereditária de parte dessa variação e da influência dessa variação na sobrevivência ou reprodução. Nesse contexto, a aptidão biológica deve ser compreendida como sucesso reprodutivo relativo em determinado ambiente, e não como uma qualidade fixa ou universal. Voltando com \citeonline{urry2021}, uma característica vantajosa em determinado ambiente pode não apresentar o mesmo valor adaptativo em outro contexto ambiental.

\subsubsection{Recombinação e Mutação}
\label{subsubsec:recombinacao-mutacao}

A recombinação genética ocorre quando materiais genéticos de diferentes progenitores são combinados na formação de descendentes. Segundo \citeonline{pierce2024}, esse processo reorganiza variantes genéticas já existentes e produz novas combinações de características, contribuindo para a diversidade entre os indivíduos.

A mutação, por sua vez, corresponde a alterações no material genético. Considerando a análise de \citeonline{futuyma2023}, os efeitos das mutações podem ser neutros, prejudiciais ou vantajosos, dependendo das condições ambientais e de sua influência sobre a sobrevivência e a reprodução. Enquanto a recombinação reorganiza informações herdadas, a mutação pode introduzir novas variações na população. Ambas contribuem para a manutenção da diversidade genética.

\subsubsection{Evolução como Processo Geracional}
\label{subsubsec:evolucao-processo-geracional}

A evolução ocorre ao longo de gerações. Em cada geração, indivíduos sobrevivem, reproduzem-se e transmitem parte de suas características. Como propõe  \citeonline{futuyma2023}, pequenas mudanças acumuladas nas frequências de características hereditárias podem alterar gradualmente a composição genética de uma população.

Esse processo depende da interação entre variabilidade, herança e seleção. Segundo \citeonline{urry2021}, a variabilidade fornece diferenças entre os indivíduos; a herança permite a transmissão dessas diferenças; e a seleção natural altera sua frequência relativa na população. A evolução não possui um objetivo final definido, pois resulta da adaptação relativa a ambientes específicos.

\subsubsection{Diversidade e Homogeneização}
\label{subsubsec:diversidade-homogeneizacao}

A diversidade genética contribui para a resiliência das populações. Segundo \citeonline{seaborn2021}, o potencial adaptativo de uma população está relacionado à sua capacidade de responder a alterações ambientais, especialmente quando existem variações hereditárias disponíveis para seleção.

A manutenção da diversidade amplia o potencial adaptativo da população, pois preserva diferentes combinações de características. A luz de  \citeonline{futuyma2023}, a redução excessiva da variabilidade pode limitar as possibilidades evolutivas de uma população ao longo das gerações. Dessa forma, a homogeneização genética pode reduzir a capacidade de adaptação diante de mudanças ambientais, doenças ou novas pressões seletivas.
